\begin{scanpage}{0060}{43}
\chaptertitle{II}{DIFFRACTION BY A DISK}
The problem of diffraction by a disk has a rigorous solution [24--26]; however, it is not suitable for numerical calculations in the quasi-optical region when the dimensions of the disk are large in comparison with the wavelength. The physical optics approach used in such cases sometimes gives erroneous results. In particular, the fringing field calculated in this approach does not satisfy the reciprocity principle.

In this Chapter a refinement of the physical optics approach is carried out. First the diffraction of a plane electromagnetic wave by a disk with normal incidence (\S\ 7--9) is investigated, and then (\S\ 10--12) diffraction by a disk with oblique incidence of a plane electromagnetic wave is investigated.
\begin{center}Normal Irradiation\end{center}
\sectiontitle{7}{The Physical Optics Approach}
Let an ideally conducting, infinitely thin disk of radius $a$ (Figure 17) be irradiated by plane wave
\end{scanpage}
